\noindent{\zihao{4}\sffamily\bfseries 摘要}
\vspace{0.6ex}

面向小零件配料与分拣，本作品设计了一套由 AMD Ryzen AI PC、Kria KV260 和六轴机械臂组成的桌面分拣系统。用户输入中文指令，指定零件类别、数量和目标托盘，系统依次完成指令解析、目标识别、轨迹规划、抓取放置与结果核验。作品以 M12、M6 螺丝和螺母为操作对象，任务执行不依赖云端推理。

AI PC 负责感知与规划：YOLO11n-seg 实例分割模型经 Quark A8W8 量化后部署在 NPU，本地 Qwen2.5-VL-3B 在 GPU 上辅助解析指令并复核结果，CPU 完成坐标变换和轨迹规划。KV260 负责执行与腕部纠偏：ARM 处理器运行作业管理和视觉定位，FPGA 可编程逻辑（PL）完成 5\ms 轨迹求值、CAN 调度、执行保护和在线修正，大螺丝的腕部图像掩码由 PL 加速。头部相机提供整盘目标位置，腕部相机在抓取前校正局部偏差，放置后再由头部相机核验结果。

四类零件验证集的掩码 mAP50 为 0.976。大模型同时运行时，NPU 检测延迟 P95 为 141\ms，相比共享 GPU 的 306\ms，减少了并发推理对持续检测的影响。CAN 总线实测 15364 个控制周期，周期间隔偏差 P50 为 0.4\us；PL 掩码处理延迟 P50 为 7.82\ms，每帧 CPU 耗时比 ARM 单核实现降低 95\%；大螺丝从悬停到对准收敛的时间 P50 为 1.06\,s。

系统还提供零重力拖动示教、关节反馈驱动的数字孪生和本地大模型中文复审。开发中使用 AMD Ross 的 HLS 技能辅助插补核检查与构建，并将代码、测试记录、部署方法和故障排查经验整理为可复用的工程资料。

\noindent\textbf{关键词：}具身智能；Ryzen AI NPU；Kria KV260；FPGA 运动控制；视觉纠偏；端侧大模型
